\chapter{Dark Pools and Blocks}\label{ch:mx:dark-pools-and-blocks}

A pension fund wants to buy two million shares of a stock that trades five million a day. Shown on a lit book, the order would move the price before a
tenth of it traded; sent to a dark pool, it waits for a seller, and every venue it touches learns a little of what it wants. Dark venues exist to let
large orders meet without announcing themselves, and they work only as well as they keep the secret. This chapter builds a dark pool and a block
venue beside a lit book in the exchange simulator, measures who fills in the dark, and works one large order four ways to price what leaks.

\section{Midpoint crossing and minimum quantities}

A dark pool (One Quant Book 1, chapter~9) displays nothing and usually crosses at the midpoint of the lit market's best bid and offer, so both sides
save half the spread against trading on the lit book: between 20.00 and 20.02, a cent a share each, USD 100 on 10\,000 shares. The price is borrowed:
the pool needs the lit quotes, and its crosses are only as current as its copy of them. \texttt{firm.exchsim} now gives a dark venue that copy: a
\emph{reference quote} control message (N) carries the lit venue's best bid and offer to the dark engine after each change, and \omterm{def:mx:order-types-and-their-uses:peg}{midpoint pegs} price off
it (\texttt{midpoint\_dark\_pool(reference=\dots)}).

\begin{definition}[Minimum execution quantity]\label{def:mx:dark-pools-and-blocks:minqty}
A \emph{minimum execution quantity}\index{minimum execution quantity} on a resting dark order is the smallest fill it accepts: an incoming order that
cannot trade at least that much with it passes it by. On an incoming \omterm{def:mx:order-types-and-their-uses:tif}{immediate-or-cancel order} it is the smallest total fill worth taking.
\end{definition}

A resting midpoint buy of 5\,000 shares with a minimum of 1\,000 ignores a 300-share sell and trades with a 1\,500-share one. The minimum is a filter on
counterparties: small orders that probe the pool, and the small orders of fast traders generally, cannot reach it. It is also a filter on fills: the
large order trades less.

\begin{definition}[Crossing network]\label{def:mx:dark-pools-and-blocks:crossing}
A \emph{crossing network}\index{crossing network} matches buy and sell orders at a price taken from elsewhere (the lit midpoint, a closing price, a
benchmark), without \omterm{def:mx:fragmentation-and-routing:discovery}{price discovery} of its own, either continuously or at scheduled times.
\end{definition}

\section{Conditional orders and block venues}

\begin{definition}[Indication of interest]\label{def:mx:dark-pools-and-blocks:ioi}
An \emph{indication of interest}\index{indication of interest} (IOI) is a non-binding message that a trader may want to trade a security, with a side
and sometimes a size or a price, sent to a counterparty or a venue to find the other side of a large trade.
\end{definition}

\begin{definition}[Conditional order]\label{def:mx:dark-pools-and-blocks:conditional}
A \emph{conditional order}\index{conditional order} is an \omterm{def:mx:dark-pools-and-blocks:ioi}{indication of interest} resting on a block venue: when it meets a contra conditional of at
least both sides' minimum sizes, the venue invites both traders to \emph{firm up}, to send a binding order, within a short timer. If both do, the block
crosses; if either does not, nothing trades.
\end{definition}

\begin{definition}[Firm-up rate]\label{def:mx:dark-pools-and-blocks:firmup}
A participant's \emph{firm-up rate}\index{firm-up rate} is the share of the invitations it receives to which it responds with a firm order. Block
venues track it and restrict participants whose rate is low, because a conditional that does not firm up has learnt that the other side exists at no cost.
\end{definition}

Conditionals let one order rest on several block venues and algorithms at once without risking a double fill, since only a firm-up commits it. Their
weakness is in the definition: an invitation tells both sides that a large contra exists. A counterparty that never meant to trade gets the information
for free, which is why the \omterm{def:mx:dark-pools-and-blocks:firmup}{firm-up rate} is policed (\cref{fig:mx:dark-pools-and-blocks:conditional}). \texttt{firm.darkpool.ConditionalBook} implements
the matching, invitations, firm-up timers and rates.

\begin{omfigure}
\begin{tikzpicture}[font=\footnotesize,box/.style={draw,rounded corners,align=center,minimum height=0.8cm}]
\node[box,fill=omDef!12] (b) at (0,1.2) {fund\\ buy 20\,000\\ min 1\,000};
\node[box,fill=omDef!12] (s) at (0,-1.2) {seller\\ sell 8\,000\\ min 2\,000};
\node[box,fill=omProp!15] (v) at (4,0) {block venue\\ match $\Rightarrow$ invite};
\node[box,fill=omThm!12] (x) at (8.3,1) {both firm up:\\ 8\,000 cross at the mid};
\node[box,fill=omExo!12] (l) at (8.3,-1) {one does not:\\ no trade, but it knows};
\draw[->] (b) -- node[above,sloped] {conditional} (v);
\draw[->] (s) -- node[below,sloped] {conditional} (v);
\draw[->] (v) -- (x);
\draw[->] (v) -- (l);
\end{tikzpicture}

\omcaption{A conditional block. Two non-binding orders meet, the venue invites both to firm up within a timer, and only two firm orders trade. A
participant who lets invitations lapse learns about the other side without trading.}\label{fig:mx:dark-pools-and-blocks:conditional}
\end{omfigure}

\omcode{code/firm/darkpool/firm_darkpool.py}{70}{89}{Matching conditionals in time priority and inviting both sides; an order is in one invitation at a time.}\label{lst:mx:dark-pools-and-blocks:invite}

\section{Who trades in the dark, and the evidence}

Zhu (2014) explained who goes dark. \omterm{def:mx:why-there-is-a-spread:informed}{Informed traders} trade in the same direction as one another, crowd on the heavy side of the market and so face a
higher risk of not being filled in the dark, where the other side must come to them; the exchange suits them better, the dark pool suits uninformed
traders, and a dark pool can concentrate information on the exchange. Comerton-Forde and Putnins (2015) found in Australian data that dark trades are
less informed than lit trades, that low levels of non-block dark trading are benign or even beneficial for informational efficiency and high levels
harmful, and no evidence that block trades in the dark impede \omterm{def:mx:fragmentation-and-routing:discovery}{price discovery}.

The chapter's market reproduces the selection. Its efficient price $P^\ast$ moves by exogenous jumps and by the permanent impact of uninformed lit
\omterm{def:mx:the-limit-order-book:orders}{market orders}, $\lambda=0.0005$ ticks a share (informed orders reveal $P^\ast$, they do not move it). Liquidity providers quote on LIT; noise and
informed takers try the dark pool first half the time and send what is left to LIT; patient uninformed traders rest \omterm{def:mx:order-types-and-their-uses:peg}{midpoint pegs} in the dark. Over four
seeds of 45 minutes, informed takers who tried the dark were filled for 6.3\% of their shares, uninformed ones for 17.5\%: the informed buy when
everyone informed buys, and the dark sellers run out. Resting in the dark gives up the spread, but its fills were followed by a move of $+0.11$ cents in the
resting trader's favour over 30 seconds, against $-0.13$ cents for lit providers, who capture 0.51 cents of spread and pay adverse selection on it.

\section{Information leakage}

\begin{definition}[Information leakage]\label{def:mx:dark-pools-and-blocks:leakage}
\emph{Information leakage}\index{information leakage} is the part of a large order's cost that comes from other traders learning of it before it is
complete, from its fills, its quotes, its conditionals or its routing, and trading ahead of it.
\end{definition}

\begin{definition}[Pinging order]\label{def:mx:dark-pools-and-blocks:ping}
A \emph{pinging order}\index{pinging order} is a small \omterm{def:mx:order-types-and-their-uses:tif}{immediate-or-cancel order} sent to a dark venue to detect hidden liquidity: a fill reveals a
resting contra order at that price, and repeated fills reveal a large one.
\end{definition}

A fund buys 20\,000 shares from ten minutes into the session over 30 minutes, four ways: \emph{lit}, a \omterm{def:mx:the-limit-order-book:orders}{market order} of 200 every 18 seconds;
\emph{dark}, one resting \omterm{def:mx:order-types-and-their-uses:peg}{midpoint peg} for the whole order; \emph{dark with a minimum}, the same peg with a minimum of 1\,000; \emph{block}, conditionals
in a block venue where natural sellers post conditionals of 5\,000 to 15\,000 shares every five minutes on average, 30\% of whom never firm up and buy
2\,000 shares on LIT instead. Whatever is left after 30 minutes is bought on LIT at once. A pinger probes the dark pool every two seconds with 100 shares,
alternating sides; when two probes in a row on one side fill within five seconds, it buys (or sells) 500 shares on LIT ahead of the order it has found.

\omcode{code/microstructure/09-dark-pools-and-blocks/python/mx_dark.py}{125}{133}{The pinger: a filled probe is noted, and two in a row trigger a trade ahead on the lit venue.}\label{lst:mx:dark-pools-and-blocks:ping}

The shortfalls are noisy, so the simulation also accounts for the two costs it knows. Each lit buy of $q$ moves $P^\ast$ by $\lambda q$ and raises the
price of every share the fund buys afterwards: summed over the fund's own lit children this is its own-impact cost, $\lambda Q/2\,(1-1/n)=4.95$ basis
points for 100 children; summed over the pinger's and the leaky sellers' trades ahead of it, when they were triggered by the fund, it is the leakage cost.
Over sixteen seeds with the same exogenous events for every strategy (\cref{fig:mx:dark-pools-and-blocks:costs}):

\begin{itemize}
\item \emph{Lit}: shortfall 3.8 basis points (standard error 1.4), drift of the mid over the window 7.3 (2.3), own impact 4.95, no leakage.
\item \emph{Dark}: the peg completes in about 360 seconds, but the pinger finds it: 88 of its probes fill against the fund, it buys 44\,000 shares ahead,
and the leakage costs 10.9 basis points (0.3). Shortfall 12.5 (0.5), drift 26.1 (2.2).
\item \emph{Dark with a minimum of 1\,000}: no probe can fill it and there is no leakage, but only 46\% of the order fills in the dark; the rest is bought
at the end, at 1.4 basis points of own impact. Shortfall 9.5 (1.6), drift 11.9 (1.6).
\item \emph{Block}: 79\% crosses in blocks, 0.75 leaky invitations per run cost 0.6 basis points. Shortfall $-2.9$ (1.9), drift $-1.7$ (2.2).
\end{itemize}

\begin{omfigure}
\begin{tikzpicture}
\begin{axis}[width=10cm,height=5.4cm,ybar,bar width=9pt,xtick={0,1,2,3},xticklabels={lit,dark,{dark, min 1\,000},block},
  ylabel={\small basis points},tick label style={font=\footnotesize},ymin=-6,ymax=30,grid=major,grid style={gray!25},
  legend columns=3,legend style={font=\footnotesize,at={(0.5,-0.2)},anchor=north,draw=none,fill=none,
  /tikz/every even column/.append style={column sep=8pt}},table/col sep=comma,enlarge x limits=0.2]
\addplot[omDef,fill=omDef!40,error bars/.cd,y dir=both,y explicit] table[x=k,y=shortfall,y error=shortfall_se]{figdata/microstructure/09-dark-pools-and-blocks/costs.csv};
\addplot[omThm,fill=omThm!40] table[x=k,y=leak]{figdata/microstructure/09-dark-pools-and-blocks/costs.csv};
\addplot[omProp,fill=omProp!40] table[x=k,y=own]{figdata/microstructure/09-dark-pools-and-blocks/costs.csv};
\legend{shortfall,leakage (accounted),own impact (accounted)}
\end{axis}
\end{tikzpicture}

\omcaption{A 20\,000-share buy worked four ways in the simulated market: shortfall against the arrival mid (mean and standard error over sixteen seeds),
and the two costs the simulation can account for exactly. Resting unprotected in the dark leaks more than trading on the lit book costs. Data:
\texttt{mx\_dark.compare}.}\label{fig:mx:dark-pools-and-blocks:costs}
\end{omfigure}

The lessons hold beyond the model's numbers. A resting dark order is visible to whoever probes for it, and the probe costs the prober 100 shares at the
midpoint; the leak can cost more than the impact the dark venue was chosen to avoid. A minimum quantity above the probe size stops the leak and costs
completion. And a midpoint fill saves the spread but not the impact when the counterparty is flow that would otherwise have sold on the lit book: taking
it out of the lit market moves the price the same way. Blocks are cheap in the simulation because their sellers bring size that would not otherwise have
traded: the cost of a trade is the imbalance it creates, wherever it is printed.

\section{Dark pools in enforcement records}

The secret is the product, and enforcement records show what happens when operators do not keep it. On August 12, 2015 ITG and its affiliate AlterNet
Securities agreed to pay USD 20.3 million, admitting that an undisclosed proprietary desk, Project Omega, had accessed live feeds of its dark pool
subscribers' orders and executions and traded against them in POSIT. On January 31, 2016 Barclays agreed to pay USD 70 million, half to the SEC and half to
the New York Attorney General, for misrepresenting how its Liquidity Profiling policed its LX pool, and Credit Suisse USD 84.3 million over its
Crossfinder pool, where it misrepresented its Alpha Scoring of subscribers' flow, accepted over 117 million sub-penny orders and alerted two
high-frequency firms to customer orders through its Crosslink service. Each is a case of \omterm{def:mx:dark-pools-and-blocks:leakage}{information leakage} by the venue itself.

\begin{dated}{2026-09}{Dark trading caps in the European Union}\label{dat:mx:dark-pools-and-blocks:svc}
The MiFIR review entered into force on 28 March 2024. It replaces the double volume cap on dark trading with a single volume cap that applies only to the
reference price waiver, from 18 months after entry into force; ESMA planned its first suspension file under the new cap for 9 October 2025.
\end{dated}

\section{Tutorial: the block that leaked}\label{tut:mx:dark-pools-and-blocks:tut}

\begin{tutorial}
\textbf{Goal.} Build a lit and a dark venue with a block layer, measure who fills in the dark, and price the leakage of a large order.
\textbf{End state:} \cref{fig:mx:dark-pools-and-blocks:costs} and the numbers of sections 3 and 4.
\begin{tutsteps}
\item \textbf{Venues.} \texttt{mx\_dark.Market(seed, strategy)}: two engines driven directly, the dark one priced by control N after each change of the
lit top.
\item \textbf{Selection.} \texttt{zhu()}: dark fill rates of informed and uninformed takers, capture and 30-second move of resting fills in the dark and on
the lit book.
\item \textbf{Leakage.} \texttt{run(seed, strategy)} for \texttt{lit}, \texttt{dark}, \texttt{dark\_min}, \texttt{block}; \texttt{compare()} over
sixteen seeds; draw with \texttt{fig\_dark.py}.
\item \textbf{Block layer.} \texttt{firm\_darkpool.ConditionalBook}: add conditionals, collect invitations, firm up, expire, and read the \omterm{def:mx:dark-pools-and-blocks:firmup}{firm-up rates}.
\end{tutsteps}
\textbf{What to change next.} Give the pinger 50-share probes and the fund a minimum of 200; let the block venue exclude participants whose \omterm{def:mx:dark-pools-and-blocks:firmup}{firm-up rate}
falls below one half, and watch the leakage.
\end{tutorial}

\section{Build: the dark venue}\label{bld:mx:dark-pools-and-blocks:darkpool}

\begin{build}
\bfield{Purpose} Dark and block liquidity for the execution algorithms of chapters 16 to 18 and 28, and the leakage measures the transaction-cost
analysis of chapter~19 reports.
\bfield{Interface} \texttt{ConditionalBook(firm\_up\_ns)}: \texttt{add}, \texttt{cancel}, \texttt{invitations}, \texttt{firm\_up}, \texttt{expire},
\texttt{firm\_up\_rate}; \texttt{shortfall\_bps}, \texttt{pre\_completion\_drift\_bps}, \texttt{counterparty\_markout}. In \texttt{firm.exchsim}: control
N in the Python, C++20 and Rust engines, \texttt{ExchangeConfig(reference, reference\_ns)}, \texttt{midpoint\_dark\_pool(reference=\dots)}.
\bfield{Rules} A conditional matches only when the smaller size meets both minimums; an order is in one invitation at a time; the block crosses only
when both firm up within the timer; a lapsed invitation counts against whoever did not respond.
\bfield{Acceptance tests} \texttt{code/firm/darkpool/tests/}: an invitation, a two-sided firm-up, a lapse and the resulting rates; shortfall, drift and
counterparty mark-out by hand. \texttt{code/firm/exchsim/tests/}: a \omterm{def:mx:order-types-and-their-uses:peg}{midpoint peg} with a minimum ignores a small order and crosses a large one at the lit
mid; the fixtures with three N messages stay byte-identical in C++20 and Rust.
\bfield{Stretch} Periodic dark auctions; venue-level participant scoring that blocks low \omterm{def:mx:dark-pools-and-blocks:firmup}{firm-up rates}; the counterparty mark-out curve of chapter~19.
\end{build}

\begin{omsources}
\item H.~Zhu, ``Do dark pools harm price discovery?'', \emph{Review of Financial Studies} 27(3), 2014.
\item C.~Comerton-Forde and T.~Putni\c{n}\v{s}, ``Dark trading and price discovery'', \emph{Journal of Financial Economics} 118(1), 2015.
\item U.S. Securities and Exchange Commission, press release 2015-164 (ITG, AlterNet Securities), 2015; press release 2016-16 (Barclays, Credit Suisse),
2016.
\item European Securities and Markets Authority, \emph{Double Volume Cap Mechanism}, web page, 2025.
\end{omsources}

\section{Exercises}

\begin{exercise}[$\star$]\label{exo:mx:dark-pools-and-blocks:1}
The lit market is 20.00 bid, 20.02 offered. What does each side of a 10\,000-share midpoint cross save against crossing the spread on the lit book?
\end{exercise}

\begin{exercise}[$\star$]\label{exo:mx:dark-pools-and-blocks:2}
A resting midpoint buy of 5\,000 with a minimum of 1\,000 meets a 300-share sell, then a 1\,500-share sell. What trades?
\end{exercise}

\begin{exercise}[$\star$]\label{exo:mx:dark-pools-and-blocks:3}
A participant received 40 invitations and firmed up 28. What is its \omterm{def:mx:dark-pools-and-blocks:firmup}{firm-up rate}, and why does the venue care?
\end{exercise}

\begin{exercise}[$\star\star$]\label{exo:mx:dark-pools-and-blocks:4}
With $\lambda=0.0005$ ticks a share on a 100.00 stock, a pinger buys 500 shares ahead when the fund has bought 5\,000 of 20\,000. How much does that one
trade cost the fund, in basis points of the whole order?
\end{exercise}

\begin{exercise}[$\star\star$]\label{exo:mx:dark-pools-and-blocks:5}
In the simulation informed takers fill 6.3\% of what they try in the dark and uninformed ones 17.5\%. Explain the gap with Zhu's argument.
\end{exercise}

\begin{exercise}[$\star\star$]\label{exo:mx:dark-pools-and-blocks:6}
Show that 100 equal lit children of a 20\,000-share order cost $\lambda Q/2\,(1-1/n)$ in own impact, and compute it.
\end{exercise}

\begin{exercise}[$\star\star\star$]\label{exo:mx:dark-pools-and-blocks:7}
\emph{Coding.} Rerun \texttt{dark\_min} with a minimum of 200 instead of 1\,000 on the sixteen seeds. Compare the leakage, the share filled in the dark,
the completion time and the shortfall with both dark strategies.
\end{exercise}

\begin{exercise}[$\star\star\star$]\label{exo:mx:dark-pools-and-blocks:8}
\emph{Find the flaw.} ``All our dark fills were at the midpoint, so our large order cost us nothing.''
\end{exercise}

\section{Problem: The Block That Leaked}

\begin{problem}[{Weekend problem --- the block that leaked}]\label{pb:mx:dark-pools-and-blocks:1}
A fund must buy 20\,000 shares within 30 minutes and asks where to send the order. Answer with the simulated market.

\medskip\noindent\textbf{Part I --- The venues.}
\begin{enumerate}
\item What does a midpoint cross save, and where does its price come from?
\item What does the dark engine need from the lit venue, and how does the simulator give it?
\item Define a \omterm{def:mx:dark-pools-and-blocks:minqty}{minimum execution quantity} and give its two effects.
\item Define a \omterm{def:mx:dark-pools-and-blocks:conditional}{conditional order} and a \omterm{def:mx:dark-pools-and-blocks:firmup}{firm-up rate}.
\item Why does a lapsed invitation leak information?
\end{enumerate}

\medskip\noindent\textbf{Part II --- Selection.}
\begin{enumerate}[resume]
\item State Zhu's argument.
\item Give the dark fill rates of informed and uninformed takers in the simulation.
\item Compare the capture and 30-second move of resting fills in the dark and on the lit book.
\item What did Comerton-Forde and Putnins find about block trades?
\end{enumerate}

\medskip\noindent\textbf{Part III --- Four ways to buy.}
\begin{enumerate}[resume]
\item Describe the pinger and the leaky sellers.
\item Derive the own-impact cost of the lit strategy.
\item How is the leakage cost accounted?
\item Give the shortfall, drift and accounted costs of the four strategies.
\item Why does the unprotected dark order cost more than the lit one?
\item What does the minimum quantity of 1\,000 fix, and what does it cost?
\end{enumerate}

\medskip\noindent\textbf{Part IV --- Venues and the verdict.}
\begin{enumerate}[resume]
\item Why are blocks cheap in the simulation, and when would they not be?
\item Summarise the ITG, Barclays and Credit Suisse cases in one sentence each.
\item What changed in the EU's cap on dark trading?
\item State the \emph{named result}: the price drift before completion of a parent order worked in the dark against in the lit market, and the leakage
cost in basis points.
\item In one sentence: where should the fund send its order?
\end{enumerate}
\end{problem}

\section{Interview questions}

\begin{interviewq}[$\star$ \iqroles{trader}]\label{iq:mx:dark-pools-and-blocks:1}
What is a dark pool, and why would an institution use one?
\end{interviewq}

\begin{interviewq}[$\star\star$ \iqroles{trader}]\label{iq:mx:dark-pools-and-blocks:2}
How would you protect a large resting dark order from being detected?
\end{interviewq}

\begin{interviewq}[$\star\star$ \iqroles{researcher}]\label{iq:mx:dark-pools-and-blocks:3}
Do dark pools harm \omterm{def:mx:fragmentation-and-routing:discovery}{price discovery}? Give the theory and the evidence.
\end{interviewq}

\begin{interviewq}[$\star\star$ \iqroles{trader, researcher}]\label{iq:mx:dark-pools-and-blocks:4}
How would you measure \omterm{def:mx:dark-pools-and-blocks:leakage}{information leakage} in your firm's executions?
\end{interviewq}

\begin{interviewq}[$\star\star$ \iqroles{developer}]\label{iq:mx:dark-pools-and-blocks:5}
Design the matching logic of a conditional block venue. What state does it keep per order and per participant?
\end{interviewq}

\begin{interviewq}[$\star\star\star$ \iqroles{trader}]\label{iq:mx:dark-pools-and-blocks:6}
Your midpoint fills look free, yet your shortfall on dark-heavy orders is higher than on lit ones. What is going on?
\end{interviewq}
